\documentclass[lettersize,journal]{IEEEtran}
% \usepackage[UTF8,fontset=windows]{ctex}
\usepackage{booktabs}
\usepackage{dsfont}
\usepackage{graphicx}
\usepackage[caption=false,font=footnotesize]{subfig}\usepackage{array}
\usepackage{textcomp}
\usepackage{verbatim}
\hyphenation{op-tical net-works semi-conduc-tor IEEE-Xplore}
\newenvironment{redsection}
  {\begingroup\color{red}
   \captionsetup{labelfont={color=red}, textfont={color=red}}}
  {\endgroup}
% \usepackage[utf8]{inputenc}
% \usepackage[T1]{fontenc}
% \usepackage{aecompl}
\usepackage{dsfont}
\usepackage{epstopdf}
\usepackage{cite}
\usepackage{amsmath,amssymb,amsfonts, amsthm}
\usepackage{algorithmicx}
\usepackage{algpseudocode}
\usepackage{algorithm}
\usepackage{xcolor}
\usepackage{float}
\usepackage[table]{xcolor}
\usepackage{booktabs}
\usepackage{array}
\usepackage{tabularx}
\newcommand{\NA}{\cellcolor{gray!18}}
\newcommand{\NoAct}{--}
% \usepackage{picinpar}
% \usepackage{babel}
\usepackage{url}
% \usepackage[latin1]{inputenc}
\usepackage{multirow}
\usepackage{color}
\usepackage{enumerate}
\usepackage{siunitx}
% \usepackage{breakurl}
\usepackage{epstopdf}
\usepackage{pbox}
\usepackage{stfloats}
\usepackage{adjustbox}
\usepackage[colorlinks=true,
            linkcolor=black,  
            citecolor=black,    
            urlcolor=blue      
]{hyperref}\newtheorem{assumption}{Assumption}
\newtheorem{definition}{Definition}
\newtheorem{theorem}{Theorem}
\newtheorem{corollary}{Corollary}
\newtheorem{proposition}{Proposition}
\newtheorem{lemma}{Lemma}
\newtheorem{remark}{Remark}
\usepackage[most]{tcolorbox}
\tcbset{highlightstyle/.style={
  colback=yellow!10, 
  colframe=red!50,   
  boxrule=0.5pt, arc=3pt, left=3pt, right=3pt, top=3pt, bottom=3pt
}}
\algrenewcommand\algorithmicrequire{\textbf{Input:}}
\algrenewcommand\algorithmicensure{\textbf{Output:}}
\usepackage{stfloats}
\def\BibTeX{{\rm B\kern-.05em{\sc i\kern-.025em b}\kern-.08em
    T\kern-.1667em\lower.7ex\hbox{E}\kern-.125emX}}
\usepackage{balance}

\begin{document}

\title{Learning Stage-wise Constraints from Expert Demonstrations}


\author{Baiyu Peng$^{1}$, Aude Billard$^{1}$%
\thanks{
This work was supported by euROBIN Grant.  } %Use only for final RAL version
\thanks{$^{1}$ The authors are with the LASA, School of Engineering, EPFL (Swiss Federal Institute of Technology in Lausanne), Lausanne 1015 Vaud, Switzerland.
        ({\tt\footnotesize baiyu.peng@epfl.ch} ; {\tt\footnotesize aude.billard@epfl.ch}). Corresponding author: Baiyu Peng.}%

}
\markboth{Journal of \LaTeX\ Class Files,~Vol.~18, No.~9, September~2024}%
{How to Use the IEEEtran \LaTeX \ Templates}

\maketitle

\begin{abstract}
Learning task constraints from demonstrations provides an explicit and reusable representation for downstream planning. However, many existing constraint-learning methods assume that the same constraint set remains active throughout an entire trajectory, which is restrictive for multi-stage robotic tasks where different constraints govern different phases and the stage cutpoints are unobserved. This paper studies stage-wise constraint learning from feasible demonstrations. Given a task-dependent feature bank, the goal is to recover, for each stage and feature, whether the feature is inactive, equality-constrained, lower-bound constrained, or upper-bound constrained, together with the corresponding target or boundary value when active. We propose Stage-Wise Constraint Learning (SWCL) to jointly segments demonstrations and recovers stage-wise constraints. For each candidate segment, SWCL evaluates local feature evidence for equality, inequality, or inactive interpretations, subgoal-progress consistency, and cross-demonstration agreement with shared stage-wise prototypes. With fixed prototypes, each demonstration is segmented exactly by dynamic programming over ordered cutpoints; the prototypes are then updated by aggregating local modes and parameters across demonstrations. Experiments on planar obstacle avoidance, slide insertion, and spherical surface inspection show that SWCL recovers more accurate cutpoints and constraint parameters than HMM-, HSMM-, and clustering-based baselines, even though the baselines are given ground-truth active constraint pairs during parameter evaluation. Planning experiments further show that the learned constraints can be reused as feasibility conditions in modified and new task instances.
\end{abstract}

\begin{IEEEkeywords}
Constraint learning, Learning from demonstration
\end{IEEEkeywords}
\section{Preliminaries and Problem Formulation}
\label{sec:preliminaries}

We consider multi-stage robotic tasks represented through a prescribed bank of
task-dependent scalar features. Let
$\mathcal{X}\subseteq\mathbb{R}^{d_x}$ denote the state space, and let
\[
\Phi=\{\phi_j:\mathcal{X}\rightarrow\mathbb{R}\}_{j=1}^{F}
\]
be the feature bank. For a state $x_t\in\mathcal{X}$, we denote the corresponding
feature value by
$u_{t,j}=\phi_j(x_t)$ and the full feature vector by
$\phi(x_t)=[u_{t,1},\dots,u_{t,F}]^\top$.
Before learning, each feature dimension is standardized over the demonstrated
data. All constraint parameters introduced below are therefore expressed in
standardized feature coordinates during learning and can subsequently be mapped
back to their original physical units.

A task consists of $K$ ordered stages. For demonstration $i$, let
\[
\tau^{(i)}
=
\{x_t^{(i)}\}_{t=1}^{T_i}
\]
denote a trajectory of length $T_i$. Its unknown stage schedule is represented by
the cutpoints
\[
B^{(i)}
=
\{b_k^{(i)}\}_{k=0}^{K},
\qquad
1=b_0^{(i)}<b_1^{(i)}<\cdots<b_K^{(i)}=T_i+1.
\]
The corresponding stage interval is
\begin{equation}
\mathcal{I}_k^{(i)}
\left(B^{(i)}\right)
=
\left\{
b_{k-1}^{(i)},\dots,b_k^{(i)}-1
\right\}.
\label{eq:stage_interval_map}
\end{equation}
The cutpoints are latent and may differ across demonstrations.

\subsection{Shared Stage-wise Constraint Model}

Each stage-feature pair is associated with a shared task-level constraint mode
\[
\bar{\gamma}_{k,j}
\in
\{\varnothing,\mathrm{eq},\mathrm{lb},\mathrm{ub}\},
\]
where $\varnothing$ denotes an inactive feature, $\mathrm{eq}$ denotes an
equality constraint, and $\mathrm{lb}$ and $\mathrm{ub}$ denote lower- and
upper-bound inequalities, respectively.

Each feature is assigned a known admissible constraint family. Let
$\mathcal{F}_{\mathrm{eq}}$ and $\mathcal{F}_{\mathrm{ineq}}$ denote the
equality-type and inequality-type feature sets. We define the admissible mode set
for feature $j$ as
\begin{equation}
\mathcal{M}_j
=
\begin{cases}
\{\varnothing,\mathrm{eq}\},
& j\in\mathcal{F}_{\mathrm{eq}},
\\
\{\varnothing,\mathrm{lb},\mathrm{ub}\},
& j\in\mathcal{F}_{\mathrm{ineq}}.
\end{cases}
\label{eq:admissible_mode_set_map}
\end{equation}

When $\bar{\gamma}_{k,j}\neq\varnothing$, the corresponding shared scalar
constraint parameter is denoted by $\bar{\eta}_{k,j}\in\mathbb{R}$. It represents
an equality target when $\bar{\gamma}_{k,j}=\mathrm{eq}$ and a boundary value
when
$\bar{\gamma}_{k,j}\in\{\mathrm{lb},\mathrm{ub}\}$.
We collect the shared variables of one stage-feature pair as
\[
\bar{\theta}_{k,j}
=
\left(
\bar{\gamma}_{k,j},
\bar{\eta}_{k,j}
\right),
\]
where $\bar{\eta}_{k,j}$ is omitted when
$\bar{\gamma}_{k,j}=\varnothing$.
The complete shared stage-wise constraint model is
\[
\bar{\Theta}
=
\{
\bar{\theta}_{k,j}
:
k=1,\dots,K,\;
j=1,\dots,F
\}.
\]

The shared mode and parameter define the corresponding feasible set
\begin{equation}
\bar{\mathcal{C}}_{k,j}
=
\begin{cases}
\mathcal{X},
& \bar{\gamma}_{k,j}=\varnothing,
\\
\{x\in\mathcal{X}:\phi_j(x)=\bar{\eta}_{k,j}\},
& \bar{\gamma}_{k,j}=\mathrm{eq},
\\
\{x\in\mathcal{X}:\phi_j(x)\geq\bar{\eta}_{k,j}\},
& \bar{\gamma}_{k,j}=\mathrm{lb},
\\
\{x\in\mathcal{X}:\phi_j(x)\leq\bar{\eta}_{k,j}\},
& \bar{\gamma}_{k,j}=\mathrm{ub}.
\end{cases}
\label{eq:shared_stage_feature_feasible_set}
\end{equation}

Known task-independent feasibility requirements, including system dynamics and
known global collision or actuator constraints, are collected in
$\mathcal{C}_{\mathrm{known}}$.

\subsection{Demonstration-specific Local Constraint Variables}

Although corresponding stages share the same task-level constraint model,
individual demonstrations may express the shared constraints imperfectly because
of execution variability, finite sampling, noise, or imperfect stage
localization. We therefore introduce a local latent constraint interpretation for
each demonstration, stage, and feature.

For demonstration $i$, let
\[
\gamma_{k,j}^{(i)}
\in
\mathcal{M}_j
\]
denote the local constraint mode of feature $j$ in stage $k$. When
$\gamma_{k,j}^{(i)}\neq\varnothing$, let
$\eta_{k,j}^{(i)}\in\mathbb{R}$ denote its local target or boundary parameter.
We write
\[
\theta_{k,j}^{(i)}
=
\left(
\gamma_{k,j}^{(i)},
\eta_{k,j}^{(i)}
\right),
\]
with $\eta_{k,j}^{(i)}$ omitted for the inactive mode, and define
\[
\Theta^{(i)}
=
\{
\theta_{k,j}^{(i)}
:
k=1,\dots,K,\;
j=1,\dots,F
\}.
\]

The variables $\Theta^{(i)}$ are auxiliary demonstration-specific latent
variables. The final task representation is the shared model $\bar{\Theta}$,
rather than the individual local interpretations.

For a candidate interval $\mathcal{I}$ in demonstration $i$, we denote the
observed values of feature $j$ by
\begin{equation}
\mathcal{U}_j^{(i)}(\mathcal{I})
=
\{
u_{t,j}^{(i)}
=
\phi_j(x_t^{(i)})
:
t\in\mathcal{I}
\}.
\label{eq:interval_feature_samples_map}
\end{equation}

We are given a set of feasible demonstrations
\[
\mathcal{D}
=
\{
\tau^{(i)}
\}_{i=1}^{N}.
\]
The demonstrations share the number and ordering of stages and the shared
constraint model $\bar{\Theta}$, while their cutpoints
$\{B^{(i)}\}$ and local constraint variables
$\{\Theta^{(i)}\}$ are unobserved.

\begin{definition}[Stage-wise constraint learning]
Given the demonstrations $\mathcal{D}$, the number of stages $K$, the feature
bank $\Phi$, and the admissible feature families
$\mathcal{F}_{\mathrm{eq}}$ and $\mathcal{F}_{\mathrm{ineq}}$, the objective is
to estimate the shared stage-wise constraint model $\bar{\Theta}$.
The demonstration-specific schedules $\{B^{(i)}\}$ and local constraint
variables $\{\Theta^{(i)}\}$ are jointly inferred as auxiliary latent variables.
\end{definition}


\section{Method}
\label{sec:method}

We formulate SWCL as regularized hierarchical maximum-a-posteriori inference.
The probabilistic part of the model contains two levels. First, the feature
observations within each stage are modeled conditionally on a
demonstration-specific local constraint interpretation. Second, the local
constraint variables are modeled as noisy realizations of a shared stage-wise
prototype. An additional subgoal-progress regularizer favors motion-coherent
candidate stages.

The resulting optimization alternates between two conditional problems.
Conditioned on the current shared prototypes, each demonstration is independently
segmented by dynamic programming while its local constraint modes and parameters
are inferred. Conditioned on the resulting local variables, the shared prototype
of every stage-feature pair is updated by MAP estimation.

\subsection{Observation Model for Local Constraints}
\label{subsec:observation_model_map}

We first specify the likelihood of the observed feature samples under each local
constraint mode. The observation-model parameters introduced below are fixed
during the alternating MAP procedure and are not part of the learned task
constraint representation.

\subsubsection{Inactive Reference Model}

For each feature $j$, let
$q_j(u)$ denote a proper background density used as the reference observation
model for the inactive mode. We use a Student-$t$ reference density,
\begin{equation}
q_j(u)
=
\frac{1}{\sigma_{0,j}}
t_{\nu_0}
\left(
\frac{u-\mu_{0,j}}{\sigma_{0,j}}
\right),
\label{eq:inactive_reference_density}
\end{equation}
where $\mu_{0,j}$ and $\sigma_{0,j}>0$ are feature-specific background
parameters and $t_{\nu_0}$ denotes the standard Student-$t$ density with
$\nu_0$ degrees of freedom. The background parameters are estimated once from
the standardized demonstration data and then held fixed during alternating
inference.

For the inactive mode,
\begin{equation}
p_j
\left(
u
\mid
\gamma=\varnothing
\right)
=
q_j(u).
\label{eq:inactive_observation_model}
\end{equation}

\subsubsection{Equality Observation Model}

For an equality interpretation with local target $\eta$, we use a centered
Student-$t$ likelihood
\begin{equation}
p_j
\left(
u
\mid
\gamma=\mathrm{eq},
\eta
\right)
=
\frac{1}{\sigma_{\mathrm{eq},j}}
t_{\nu}
\left(
\frac{u-\eta}{\sigma_{\mathrm{eq},j}}
\right),
\label{eq:equality_student_t_likelihood}
\end{equation}
where $\sigma_{\mathrm{eq},j}>0$ is the observation scale of feature $j$ under
an equality interpretation and $\nu$ is the Student-$t$ degrees of freedom.
The scale parameter controls the expected demonstration variability around a
fixed equality target; it is an observation-model parameter rather than a task
constraint parameter.

\subsubsection{Inequality Observation Model}

For an inequality mode
$m\in\{\mathrm{lb},\mathrm{ub}\}$ and candidate boundary $\eta$, define the
signed slack
\begin{equation}
y_m(u;\eta)
=
\begin{cases}
u-\eta,
& m=\mathrm{lb},
\\
\eta-u,
& m=\mathrm{ub}.
\end{cases}
\label{eq:map_inequality_slack}
\end{equation}
Thus, $y_m(u;\eta)\geq0$ indicates the feasible side of the candidate boundary.

We use the soft half-Student-$t$ density
\begin{equation}
p_j
\left(
u
\mid
\gamma=m,
\eta
\right)
=
\frac{1}{Z_{\nu,\kappa}}
\begin{cases}
\dfrac{2}{\sigma_{\mathrm{ineq},j}}
t_\nu
\left(
\dfrac{y_m(u;\eta)}
{\sigma_{\mathrm{ineq},j}}
\right),
& y_m(u;\eta)\geq0,
\\[1em]
\dfrac{2}{\sigma_{\mathrm{ineq},j}}
t_\nu(0)
\exp
\left(
\dfrac{y_m(u;\eta)}
{\kappa\sigma_{\mathrm{ineq},j}}
\right),
& y_m(u;\eta)<0,
\end{cases}
\label{eq:map_soft_half_t_density}
\end{equation}
where $\sigma_{\mathrm{ineq},j}>0$ controls feasible-side dispersion,
$\kappa>0$ controls the softness of infeasible-side violations, and
\begin{equation}
Z_{\nu,\kappa}
=
1+2\kappa t_\nu(0)
\label{eq:map_half_t_normalizer}
\end{equation}
is the normalization constant.

Different from the original quantile-based local fitting rule, the inequality
boundary $\eta$ is now a latent parameter optimized directly under the
likelihood and hierarchical prior.

\subsubsection{Factorized Segment Likelihood}

For a candidate interval $\mathcal{I}$ and local constraint variables
$\theta_j=(\gamma_j,\eta_j)$, we assume conditional independence of the feature
observations across time and feature dimensions. The interval likelihood is
therefore
\begin{equation}
p
\left(
\{\mathcal{U}_j(\mathcal{I})\}_{j=1}^{F}
\mid
\{\theta_j\}_{j=1}^{F}
\right)
=
\prod_{j=1}^{F}
\prod_{t\in\mathcal{I}}
p_j
\left(
u_{t,j}
\mid
\gamma_j,\eta_j
\right),
\label{eq:factorized_segment_likelihood}
\end{equation}
where $\eta_j$ is omitted when $\gamma_j=\varnothing$.
This factorization defines the observation model used by SWCL.


\subsection{Hierarchical Model for Shared and Local Constraints}
\label{subsec:hierarchical_constraint_model}

The shared stage-wise prototype influences the local constraint interpretation
through a hierarchical prior.

Let
\[
M_j
=
|\mathcal{M}_j|
\]
denote the number of admissible modes for feature $j$. Given a shared mode
$\bar{\gamma}_{k,j}$, the local mode follows the symmetric categorical deviation
model
\begin{equation}
p
\left(
\gamma_{k,j}^{(i)}=m
\mid
\bar{\gamma}_{k,j}=\bar{m}
\right)
=
\begin{cases}
1-\epsilon,
& m=\bar{m},
\\[0.3em]
\dfrac{\epsilon}{M_j-1},
& m\neq\bar{m},
\end{cases}
\label{eq:hierarchical_mode_model}
\end{equation}
where $0<\epsilon<1$ controls the expected probability of a
demonstration-specific mode deviation.

For active local modes, we also require a prior on the local parameter.
Let
\begin{equation}
r_j(\eta)
=
\mathcal{N}
\left(
\eta;
\mu_{\eta,j},
s_{\eta,j}^{2}
\right)
\label{eq:broad_parameter_reference_prior}
\end{equation}
denote a broad proper reference prior on feature-$j$ constraint parameters.
Its parameters $\mu_{\eta,j}$ and $s_{\eta,j}$ are fixed.

Conditioned on a matching active shared mode, the local parameter is modeled by
a Laplace distribution centered at the shared parameter:
\begin{equation}
p
\left(
\eta_{k,j}^{(i)}
\mid
\gamma_{k,j}^{(i)}
=
\bar{\gamma}_{k,j}
\neq\varnothing,
\bar{\eta}_{k,j}
\right)
=
\operatorname{Lap}
\left(
\eta_{k,j}^{(i)};
\bar{\eta}_{k,j},
b_j
\right),
\label{eq:hierarchical_parameter_laplace}
\end{equation}
where $b_j>0$ controls the expected cross-demonstration variation of the local
parameter around the shared parameter.

When the local mode is active but differs from the shared mode, its parameter is
instead assigned the broad reference prior $r_j$:
\begin{equation}
p
\left(
\eta_{k,j}^{(i)}
\mid
\gamma_{k,j}^{(i)}\neq\bar{\gamma}_{k,j},
\gamma_{k,j}^{(i)}\neq\varnothing
\right)
=
r_j
\left(
\eta_{k,j}^{(i)}
\right).
\label{eq:mismatched_parameter_prior}
\end{equation}

Combining the mode and parameter terms, the complete conditional prior of a local
stage-feature variable is
\begin{equation}
p
\left(
\theta_{k,j}^{(i)}
\mid
\bar{\theta}_{k,j}
\right)
=
p
\left(
\gamma_{k,j}^{(i)}
\mid
\bar{\gamma}_{k,j}
\right)
\times
\begin{cases}
1,
& \gamma_{k,j}^{(i)}=\varnothing,
\\
\operatorname{Lap}
\left(
\eta_{k,j}^{(i)};
\bar{\eta}_{k,j},b_j
\right),
&
\gamma_{k,j}^{(i)}
=
\bar{\gamma}_{k,j}
\neq\varnothing,
\\
r_j
\left(
\eta_{k,j}^{(i)}
\right),
&
\gamma_{k,j}^{(i)}
\neq\bar{\gamma}_{k,j},
\;
\gamma_{k,j}^{(i)}
\neq\varnothing.
\end{cases}
\label{eq:complete_local_hierarchical_prior}
\end{equation}

Finally, we place a prior on each shared prototype. Let
$\pi_{j,m}>0$ denote the prior probability of shared mode $m\in\mathcal{M}_j$,
with
$\sum_{m\in\mathcal{M}_j}\pi_{j,m}=1$.
We use
\begin{equation}
p
\left(
\bar{\theta}_{k,j}
\right)
=
\pi_{j,\bar{\gamma}_{k,j}}
\times
\begin{cases}
1,
& \bar{\gamma}_{k,j}=\varnothing,
\\
r_j
\left(
\bar{\eta}_{k,j}
\right),
& \bar{\gamma}_{k,j}\neq\varnothing.
\end{cases}
\label{eq:shared_prototype_prior}
\end{equation}
A uniform mode prior can be used when no prior preference among admissible modes
is available.


\subsection{Regularized Joint MAP Objective}
\label{subsec:regularized_joint_map}

For a candidate segmentation $B^{(i)}$, let
$\mathcal{U}_{k,j}^{(i)}(B^{(i)})$ denote the feature-$j$ observations belonging
to stage $k$.
The negative log observation likelihood of one local stage-feature variable is
\begin{equation}
\mathcal{L}_{\mathrm{obs}}
\left(
\mathcal{U}_{k,j}^{(i)};
\theta_{k,j}^{(i)}
\right)
=
-
\sum_{t\in\mathcal{I}_k^{(i)}}
\log
p_j
\left(
u_{t,j}^{(i)}
\mid
\gamma_{k,j}^{(i)},
\eta_{k,j}^{(i)}
\right).
\label{eq:local_observation_nll}
\end{equation}

The hierarchical coupling cost is
\begin{equation}
\mathcal{L}_{\mathrm{hier}}
\left(
\theta_{k,j}^{(i)};
\bar{\theta}_{k,j}
\right)
=
-
\log
p
\left(
\theta_{k,j}^{(i)}
\mid
\bar{\theta}_{k,j}
\right),
\label{eq:local_hierarchical_nll}
\end{equation}
and the shared-prototype prior cost is
\begin{equation}
\mathcal{L}_{\mathrm{shared}}
\left(
\bar{\theta}_{k,j}
\right)
=
-
\log
p
\left(
\bar{\theta}_{k,j}
\right).
\label{eq:shared_prior_nll}
\end{equation}

The probabilistic component of the objective is therefore the negative log joint
posterior, up to an additive constant. We additionally retain the subgoal-progress
term as a structural regularizer. The complete optimization objective is
\begin{equation}
\begin{aligned}
\mathcal{J}
\left(
\bar{\Theta},
\{B^{(i)},\Theta^{(i)}\}_{i=1}^{N}
\right)
={}&
\sum_{i=1}^{N}
\sum_{k=1}^{K}
\sum_{j=1}^{F}
\Big[
\mathcal{L}_{\mathrm{obs}}
\left(
\mathcal{U}_{k,j}^{(i)};
\theta_{k,j}^{(i)}
\right)
\\
&\qquad\qquad+
\mathcal{L}_{\mathrm{hier}}
\left(
\theta_{k,j}^{(i)};
\bar{\theta}_{k,j}
\right)
\Big]
\\
&+
\sum_{k=1}^{K}
\sum_{j=1}^{F}
\mathcal{L}_{\mathrm{shared}}
\left(
\bar{\theta}_{k,j}
\right)
\\
&+
\lambda_{\mathrm{prog}}
\sum_{i=1}^{N}
\sum_{k=1}^{K}
C_{\mathrm{prog}}
\left(
\mathcal{I}_k^{(i)}
\right).
\end{aligned}
\label{eq:regularized_map_objective}
\end{equation}

When $\lambda_{\mathrm{prog}}=0$, the objective reduces to the negative log joint
posterior of the hierarchical constraint model. For
$\lambda_{\mathrm{prog}}>0$, SWCL performs regularized MAP estimation.


\subsection{Local MAP Model Selection for a Candidate Interval}
\label{subsec:local_map_candidate}

We now describe how the objective in
\eqref{eq:regularized_map_objective} produces an additive stage cost suitable for
dynamic programming.

Consider candidate interval $\mathcal{I}=[s,e]$ in demonstration $i$, and suppose
that this interval is assigned to stage $k$. Conditioned on the current shared
prototype $\bar{\theta}_{k,j}$, each admissible local mode
$m\in\mathcal{M}_j$ defines a candidate posterior score.

For the inactive mode,
\begin{equation}
S_{k,j}^{(i)}
\left(
\mathcal{I},\varnothing;
\bar{\theta}_{k,j}
\right)
=
-
\sum_{t\in\mathcal{I}}
\log q_j
\left(
u_{t,j}^{(i)}
\right)
-
\log
p
\left(
\gamma_{k,j}^{(i)}=\varnothing
\mid
\bar{\gamma}_{k,j}
\right).
\label{eq:inactive_candidate_map_score}
\end{equation}

For an active candidate mode
$m\in\mathcal{M}_j\setminus\{\varnothing\}$, define
\begin{equation}
\begin{aligned}
S_{k,j}^{(i)}
\left(
\mathcal{I},m;
\bar{\theta}_{k,j}
\right)
=
\min_{\eta\in\mathbb{R}}
\Bigg[
&
-
\sum_{t\in\mathcal{I}}
\log
p_j
\left(
u_{t,j}^{(i)}
\mid
\gamma=m,\eta
\right)
\\
&
-
\log
p
\left(
\gamma_{k,j}^{(i)}=m
\mid
\bar{\gamma}_{k,j}
\right)
\\
&
-
\log
p
\left(
\eta
\mid
\gamma_{k,j}^{(i)}=m,
\bar{\theta}_{k,j}
\right)
\Bigg].
\end{aligned}
\label{eq:active_candidate_map_score}
\end{equation}

Let
$\eta_{k,j}^{(i)\star}(\mathcal{I},m)$ denote the minimizing parameter in
\eqref{eq:active_candidate_map_score}. Because the constraint parameter is
scalar, this inner optimization is one-dimensional.

The local mode assigned to the candidate interval is
\begin{equation}
\gamma_{k,j}^{(i)\star}(\mathcal{I})
=
\arg\min_{m\in\mathcal{M}_j}
S_{k,j}^{(i)}
\left(
\mathcal{I},m;
\bar{\theta}_{k,j}
\right),
\label{eq:local_mode_map_selection}
\end{equation}
and its corresponding local parameter is
\begin{equation}
\eta_{k,j}^{(i)\star}(\mathcal{I})
=
\eta_{k,j}^{(i)\star}
\left(
\mathcal{I},
\gamma_{k,j}^{(i)\star}(\mathcal{I})
\right)
\label{eq:local_parameter_map_selection}
\end{equation}
when the selected mode is active.

The profiled local constraint cost is
\begin{equation}
C_{\mathrm{con},k,j}^{(i)}
\left(
\mathcal{I};
\bar{\theta}_{k,j}
\right)
=
\min_{m\in\mathcal{M}_j}
S_{k,j}^{(i)}
\left(
\mathcal{I},m;
\bar{\theta}_{k,j}
\right).
\label{eq:profiled_local_constraint_cost}
\end{equation}

Thus, unlike a two-step procedure that first estimates a local constraint
independently and subsequently penalizes its deviation from the shared
prototype, the local mode and parameter in
\eqref{eq:local_mode_map_selection}--\eqref{eq:local_parameter_map_selection}
are estimated while already conditioning on the shared stage-wise prototype.


\subsection{Subgoal-progress Regularization}
\label{subsec:subgoal_progress_map}

Constraint likelihood alone may assign a low cost to a long interval containing
multiple distinct motions. We therefore retain an auxiliary motion-coherence
regularizer based on progress toward the candidate segment endpoint.

For candidate interval
$\mathcal{I}=[s,e]$, define its endpoint-induced subgoal as
\begin{equation}
g(\mathcal{I})
=
x_e.
\label{eq:candidate_endpoint_subgoal_map}
\end{equation}
Let
\begin{equation}
d_t(\mathcal{I})
=
\left\|
x_t-g(\mathcal{I})
\right\|_2.
\label{eq:candidate_subgoal_distance_map}
\end{equation}
The progress regularizer is
\begin{equation}
C_{\mathrm{prog}}(\mathcal{I})
=
\sum_{t=s}^{e-1}
\operatorname{softplus}_{\beta}
\left(
d_{t+1}(\mathcal{I})
-
d_t(\mathcal{I})
\right),
\label{eq:progress_regularizer_map}
\end{equation}
where
$\operatorname{softplus}_{\beta}(z)
=
\beta^{-1}\log(1+\exp(\beta z))$.
Motion toward the endpoint produces a near-zero penalty, whereas motion away
from the endpoint increases the cost.

This term is not interpreted as a task constraint or as part of the feature
observation likelihood. It serves only as a structural regularizer on the
segmentation.


\subsection{Dynamic Programming Conditioned on Shared Prototypes}
\label{subsec:map_dynamic_programming}

For candidate interval $\mathcal{I}$ assigned to stage $k$ in demonstration $i$,
the complete stage cost conditioned on the current shared prototype is
\begin{equation}
C_k^{(i)}
\left(
\mathcal{I};
\bar{\Theta}
\right)
=
\sum_{j=1}^{F}
C_{\mathrm{con},k,j}^{(i)}
\left(
\mathcal{I};
\bar{\theta}_{k,j}
\right)
+
\lambda_{\mathrm{prog}}
C_{\mathrm{prog}}^{(i)}
\left(
\mathcal{I}
\right).
\label{eq:map_stage_cost}
\end{equation}

The shared-prototype prior terms
$\mathcal{L}_{\mathrm{shared}}(\bar{\theta}_{k,j})$ do not appear in
\eqref{eq:map_stage_cost}, because they are constant with respect to the
segmentation when $\bar{\Theta}$ is fixed.

For a candidate schedule $B^{(i)}$, the conditional objective of demonstration
$i$ is additive over stages:
\begin{equation}
\mathcal{J}_i
\left(
B^{(i)},\Theta^{(i)}
\mid
\bar{\Theta}
\right)
=
\sum_{k=1}^{K}
C_k^{(i)}
\left(
\mathcal{I}_k^{(i)};
\bar{\Theta}
\right).
\label{eq:conditional_demo_map_objective}
\end{equation}

Let $F_k^{(i)}(t)$ denote the minimum conditional cost of assigning the prefix
$\{x_1^{(i)},\dots,x_t^{(i)}\}$ to the first $k$ stages. The dynamic-programming
recurrence is
\begin{equation}
F_k^{(i)}(t)
=
\min_{k\leq s\leq t}
\left[
F_{k-1}^{(i)}(s-1)
+
C_k^{(i)}
\left(
[s,t];
\bar{\Theta}
\right)
\right].
\label{eq:map_dp_recurrence}
\end{equation}
The initialization is
\[
F_0^{(i)}(0)=0,
\qquad
F_0^{(i)}(t)=+\infty
\quad
\text{for }t>0.
\]

The optimal conditional cost is $F_K^{(i)}(T_i)$. Backtracking through the
minimizing start indices in \eqref{eq:map_dp_recurrence} recovers the cutpoints
$B^{(i)}$. Because every candidate stage cost stores the minimizing local mode
and parameter from
\eqref{eq:local_mode_map_selection}--\eqref{eq:local_parameter_map_selection},
the same backtracking step also recovers the corresponding local constraint
variables $\Theta^{(i)}$.

Therefore, conditioned on $\bar{\Theta}$, the DP step jointly solves
\begin{equation}
\left(
B^{(i)\star},
\Theta^{(i)\star}
\right)
=
\arg\min_{B^{(i)},\Theta^{(i)}}
\mathcal{J}_i
\left(
B^{(i)},\Theta^{(i)}
\mid
\bar{\Theta}
\right),
\label{eq:conditional_local_map_problem}
\end{equation}
where the combinatorial search over ordered stage boundaries is solved exactly
by dynamic programming and each active scalar local parameter is obtained from
the one-dimensional optimization in
\eqref{eq:active_candidate_map_score}.


\subsection{MAP Update of the Shared Constraint Prototypes}
\label{subsec:shared_prototype_map_update}

After all demonstrations have been segmented, the cutpoints
$\{B^{(i)}\}$ and local variables $\{\Theta^{(i)}\}$ are held fixed.
The observation-likelihood and progress terms in
\eqref{eq:regularized_map_objective} are then constant with respect to the
shared prototypes. Consequently, every shared stage-feature pair
$\bar{\theta}_{k,j}$ can be updated independently.

For a candidate shared mode
$m\in\mathcal{M}_j$, define the prototype objective
\begin{equation}
\begin{aligned}
G_{k,j}(m,\bar{\eta})
=
&
-
\log
p
\left(
\bar{\theta}_{k,j}
=
(m,\bar{\eta})
\right)
\\
&
-
\sum_{i=1}^{N}
\log
p
\left(
\theta_{k,j}^{(i)}
\mid
\bar{\theta}_{k,j}
=
(m,\bar{\eta})
\right),
\end{aligned}
\label{eq:shared_candidate_mode_objective}
\end{equation}
where $\bar{\eta}$ is omitted when $m=\varnothing$.

For the inactive candidate mode,
\begin{equation}
G_{k,j}^{\star}(\varnothing)
=
-
\log \pi_{j,\varnothing}
-
\sum_{i=1}^{N}
\log
p
\left(
\theta_{k,j}^{(i)}
\mid
\bar{\gamma}_{k,j}=\varnothing
\right).
\label{eq:shared_inactive_score}
\end{equation}

For each active candidate mode
$m\in\mathcal{M}_j\setminus\{\varnothing\}$, the optimal shared parameter is
\begin{equation}
\bar{\eta}_{k,j}^{\star}(m)
=
\arg\min_{\bar{\eta}\in\mathbb{R}}
G_{k,j}(m,\bar{\eta}),
\label{eq:shared_parameter_map_update}
\end{equation}
and the corresponding optimized mode score is
\begin{equation}
G_{k,j}^{\star}(m)
=
G_{k,j}
\left(
m,
\bar{\eta}_{k,j}^{\star}(m)
\right).
\label{eq:optimized_shared_mode_score}
\end{equation}

The updated shared mode is then
\begin{equation}
\bar{\gamma}_{k,j}^{\star}
=
\arg\min_{m\in\mathcal{M}_j}
G_{k,j}^{\star}(m),
\label{eq:shared_mode_map_update}
\end{equation}
and, when the selected mode is active,
\begin{equation}
\bar{\eta}_{k,j}^{\star}
=
\bar{\eta}_{k,j}^{\star}
\left(
\bar{\gamma}_{k,j}^{\star}
\right).
\label{eq:shared_parameter_selected_update}
\end{equation}

Under the Laplace coupling in
\eqref{eq:hierarchical_parameter_laplace}, only demonstrations whose local mode
matches the candidate shared mode contribute a prototype-dependent absolute
deviation term. Ignoring the weak reference prior $r_j$ on the shared
parameter, the minimizer in
\eqref{eq:shared_parameter_map_update} reduces to the median of the
mode-matched local parameter estimates. The proper reference prior in
\eqref{eq:shared_prototype_prior} turns this into a regularized one-dimensional
MAP estimate.

Importantly, the shared mode is not determined by majority vote alone.
For each candidate mode, both local-mode agreement and the consistency of its
associated local parameter estimates contribute to
$G_{k,j}^{\star}(m)$. The selected shared mode and parameter therefore form a
joint MAP prototype update.


\subsection{Deterministic Initialization and Alternating MAP Inference}
\label{subsec:alternating_map_inference}

The conditional local MAP problem requires an initial shared prototype.
We obtain it through a deterministic uncoupled initialization.

Before cross-demonstration coupling is introduced, every candidate interval is
evaluated using only the observation likelihood, a fixed local mode prior
$\pi_{j,m}$, and the broad parameter prior $r_j$. Specifically, for an active
mode $m$ we define
\begin{equation}
\begin{aligned}
S_{j}^{\mathrm{init}}
(\mathcal{I},m)
=
\min_{\eta}
\Bigg[
&
-
\sum_{t\in\mathcal{I}}
\log
p_j
\left(
u_{t,j}
\mid
m,\eta
\right)
\\
&
-
\log \pi_{j,m}
-
\log r_j(\eta)
\Bigg],
\end{aligned}
\label{eq:initial_local_map_score}
\end{equation}
while for the inactive mode
\begin{equation}
S_{j}^{\mathrm{init}}
(\mathcal{I},\varnothing)
=
-
\sum_{t\in\mathcal{I}}
\log q_j(u_{t,j})
-
\log\pi_{j,\varnothing}.
\label{eq:initial_inactive_map_score}
\end{equation}

These scores are combined with the progress regularizer and used in the same DP
recurrence as \eqref{eq:map_dp_recurrence}. The resulting initial segmentations
and local variables are then used in
\eqref{eq:shared_parameter_map_update}--\eqref{eq:shared_mode_map_update} to
construct the first shared prototype $\bar{\Theta}^{(0)}$.

Starting from $\bar{\Theta}^{(0)}$, SWCL alternates the following two updates:

\begin{enumerate}
    \item \textbf{Conditional local MAP and segmentation.}
    With $\bar{\Theta}^{(q)}$ fixed, independently solve
    \eqref{eq:conditional_local_map_problem} for every demonstration using
    dynamic programming, producing
    $\{B^{(i,q+1)},\Theta^{(i,q+1)}\}_{i=1}^{N}$.

    \item \textbf{Shared prototype MAP update.}
    With the recovered local variables fixed, update every stage-feature
    prototype independently using
    \eqref{eq:shared_parameter_map_update}--\eqref{eq:shared_mode_map_update},
    producing $\bar{\Theta}^{(q+1)}$.
\end{enumerate}

The alternating procedure terminates when the cutpoints and shared prototypes
remain unchanged, when the relative reduction in
$\mathcal{J}$ falls below a tolerance, or when a maximum number of iterations is
reached.

Because the hierarchical and observation-model hyperparameters are fixed after
initialization, both alternating updates optimize the same objective
\eqref{eq:regularized_map_objective}. If the scalar inner parameter
optimizations are solved exactly, the DP step minimizes the objective with
respect to all demonstration-specific cutpoints and local constraint variables
conditioned on the current shared prototypes, while the prototype step minimizes
it with respect to the shared variables conditioned on the current local
variables. Hence,
\begin{equation}
\mathcal{J}^{(q+1)}
\leq
\mathcal{J}^{(q)},
\label{eq:map_monotonicity}
\end{equation}
although the non-convex mixed discrete-continuous problem does not in general
guarantee convergence to the global optimum.


\begin{algorithm}[t]
\caption{Regularized Hierarchical MAP Stage-wise Constraint Learning}
\label{alg:swcl_map}
\begin{algorithmic}[1]
\Require demonstrations $\mathcal{D}$, number of stages $K$,
feature bank $\Phi$, admissible feature families
$\mathcal{F}_{\mathrm{eq}}$ and $\mathcal{F}_{\mathrm{ineq}}$,
observation-model parameters, hierarchical-model parameters,
and $\lambda_{\mathrm{prog}}$
\Ensure shared constraint model $\hat{\bar{\Theta}}$,
cutpoints $\{\hat B^{(i)}\}$,
and local latent constraints $\{\hat\Theta^{(i)}\}$

\State Extract and standardize all feature trajectories
\State Fit the fixed inactive reference densities $\{q_j\}_{j=1}^{F}$
\State Precompute the subgoal-progress costs for all candidate intervals

\Statex
\State \textbf{Deterministic initialization}
\ForAll{demonstrations $i=1,\dots,N$}
    \State Compute the uncoupled candidate scores using
    \eqref{eq:initial_local_map_score} and
    \eqref{eq:initial_inactive_map_score}
    \State Solve the DP recurrence and recover
    $B^{(i,0)}$ and $\Theta^{(i,0)}$
\EndFor
\State Update the shared prototype $\bar{\Theta}^{(0)}$ using
\eqref{eq:shared_parameter_map_update}--\eqref{eq:shared_mode_map_update}

\Statex
\State $q\leftarrow0$
\Repeat
    \ForAll{demonstrations $i=1,\dots,N$}
        \ForAll{candidate stage intervals}
            \State Compute local MAP mode and parameter scores using
            \eqref{eq:inactive_candidate_map_score}--\eqref{eq:local_mode_map_selection}
        \EndFor
        \State Solve \eqref{eq:map_dp_recurrence}
        \State Backtrack to recover
        $B^{(i,q+1)}$ and $\Theta^{(i,q+1)}$
    \EndFor

    \For{$k=1,\dots,K$}
        \For{$j=1,\dots,F$}
            \State Enumerate all $m\in\mathcal{M}_j$
            \State For each active $m$, solve
            \eqref{eq:shared_parameter_map_update}
            \State Select the shared mode using
            \eqref{eq:shared_mode_map_update}
        \EndFor
    \EndFor

    \State $q\leftarrow q+1$
\Until{convergence or maximum number of iterations}
\State \Return
$\hat{\bar{\Theta}}=\bar{\Theta}^{(q)}$,
$\{\hat B^{(i)}\}$,
and $\{\hat\Theta^{(i)}\}$
\end{algorithmic}
\end{algorithm}


\subsection{Implementation Considerations}
\label{subsec:map_implementation}

The reformulated local MAP fit depends on the current shared prototype through
the hierarchical prior. Therefore, unlike an uncoupled local-fitting procedure,
the final local mode and parameter of a candidate interval cannot in general be
computed once and reused in every outer iteration.

Quantities independent of the shared prototypes can still be cached. These
include the standardized feature trajectories, candidate interval indices,
inactive-reference likelihood contributions, and subgoal-progress costs.
Observation-model statistics that accelerate the one-dimensional active
parameter optimizations may also be precomputed. At each outer iteration, only
the prototype-dependent local MAP fits and the subsequent DP stage costs need to
be updated.

\subsection{Discussion}
\label{subsec:map_discussion}

The hierarchical MAP formulation provides a unified interpretation of the main
components of SWCL. In particular, local feature evidence is represented by an
explicit observation likelihood, while agreement between demonstration-specific
constraint interpretations and the shared stage-wise prototypes is represented
through a hierarchical prior. Consequently, local constraint fitting and
cross-demonstration consensus are no longer treated as two separate operations:
the local mode and parameter of a candidate interval are inferred while already
conditioning on the current shared prototype. This allows the data likelihood
and the shared task structure to jointly determine the local interpretation.

A second advantage is that the formulation preserves demonstration-specific
constraint variables while maintaining a shared task-level representation. A
strict shared model would require all corresponding stages to be explained
directly by exactly the same mode and parameter. In contrast, the hierarchical
model permits a local constraint estimate to deviate from the shared prototype
when the evidence in an individual demonstration supports such a deviation.
This flexibility is useful in the presence of execution variability, observation
noise, imperfect segmentation, or individual demonstrations that express a
constraint less clearly than others. At the same time, the hierarchical prior
discourages arbitrary local deviations and allows information from corresponding
stages in other demonstrations to resolve locally ambiguous interpretations.

The probabilistic formulation also makes several quantities that would otherwise
appear as independent cost weights more interpretable. For example, the
local-to-shared parameter coupling is controlled by the scale $b_j$ of the
hierarchical parameter distribution, while the strength of mode agreement is
controlled by the deviation probability $\epsilon$. These quantities have a
direct statistical interpretation as expected cross-demonstration variability.
Similarly, the observation scales
$\sigma_{\mathrm{eq},j}$ and $\sigma_{\mathrm{ineq},j}$ describe the variability
of feature observations under the corresponding active constraint models.
This does not eliminate hyperparameters, but it separates different sources of
variability and gives their relative strengths an explicit modeling
interpretation.

The formulation remains compatible with exact dynamic programming for the
segmentation subproblem. Conditioned on the current shared prototypes, the
demonstrations are independent and the objective is additive over ordered
candidate stages. The local mode and parameter of each candidate interval can
therefore be profiled out before evaluating its stage cost, after which dynamic
programming exactly searches over all admissible ordered cutpoint sequences.
The subsequent shared-prototype update is also separable over stage-feature
pairs. Hence, although the complete problem is mixed discrete-continuous and
non-convex, both alternating blocks have well-defined optimization objectives.
With fixed model hyperparameters and exact inner optimizations, the complete
regularized objective in \eqref{eq:regularized_map_objective} is non-increasing
across alternating iterations.

The probabilistic interpretation, however, introduces additional modeling
assumptions. Most importantly, an inactive feature does not have a unique
generative distribution. The absence of a task constraint does not imply that
the corresponding feature values are uniformly distributed or highly variable:
unconstrained behavior may still be concentrated because of the demonstrator's
policy, system dynamics, other task constraints, or the geometry of the task.
The reference density $q_j$ therefore acts as an explicit background model
against which active constraint hypotheses are compared, and its specification
can affect active--inactive model selection. This issue is intrinsic to
constraint inference from feasible demonstrations and is not removed by the MAP
formulation.

Relatedly, the observation scales
$\sigma_{\mathrm{eq},j}$ and $\sigma_{\mathrm{ineq},j}$ must distinguish typical
demonstration variability from the structural signature of an active
constraint. Allowing these scales to vary freely for every candidate interval
would make the active models overly flexible: a poorly fitting interval could
be explained by simply increasing its observation scale. We therefore treat
these quantities as feature-level observation-model parameters rather than
stage-specific constraint parameters. They may be calibrated from the
demonstrations or estimated at a higher level of the model, but their values
still encode assumptions about the expected variability under an active
constraint.

Another limitation is that the likelihood in
\eqref{eq:factorized_segment_likelihood} uses a conditional-independence
approximation across time and feature dimensions. Feature values computed from
the same robot state can be correlated, and neighboring trajectory samples are
generally temporally dependent. Modeling the full joint feature trajectory
distribution would require substantially more complex observation models and
would reduce the simplicity of the segment-wise dynamic-programming
formulation. The factorized likelihood should therefore be interpreted as a
tractable feature-level observation model rather than a complete generative
model of demonstrated trajectories.

The MAP formulation also increases computational cost relative to a fully
precomputed local-fitting procedure. Because the optimal local parameter of a
candidate interval depends on the current shared prototype through the
hierarchical prior, the final local MAP fit must in general be recomputed when
the prototype changes. Prototype-independent quantities, including feature
trajectories, progress costs, inactive likelihood contributions, and useful
interval statistics, can still be cached, but the complete candidate-interval
fit can no longer be precomputed once and reused throughout all outer
iterations.

Finally, the subgoal-progress term remains an auxiliary structural
regularizer rather than part of the probabilistic constraint model. It provides
additional information for distinguishing motion-coherent stages when feature
likelihood alone is insufficient, but its weight
$\lambda_{\mathrm{prog}}$ is not determined by the hierarchical probability
model. Accordingly, the complete method is more precisely described as
regularized hierarchical MAP inference rather than pure MAP estimation.
This separation is intentional: the probabilistic model describes how local and
shared constraints explain feature observations, whereas the progress term
encodes an additional assumption about the temporal coherence of a valid stage.

Overall, the reformulation trades implementation simplicity for a more unified
statistical treatment of local constraint inference and cross-demonstration
sharing. It provides a consistent objective for segmentation, local constraint
estimation, and shared prototype updates, while retaining the exact
dynamic-programming structure of the segmentation step. Its main remaining
modeling choices concern the inactive reference distribution, observation-scale
calibration, and the strength of the hierarchical coupling, all of which reflect
information that cannot in general be determined uniquely from feasible
demonstrations alone.

\bibliographystyle{IEEEtran}
\bibliography{ref}
% \begin{thebibliography}{1}
% \bibliographystyle{IEEEtran}


\end{document}Unable to help with that. Please provide the text of the paragraph for further assistance.
